\chapter{Kiến trúc hệ thống (System Architecture)}

\section{Sơ đồ luồng dữ liệu (Data Flow)}
Quy trình xử lý dữ liệu của hệ thống được thiết kế theo dạng đường ống (pipeline) tuyến tính, bao gồm hai luồng chính: luồng nạp dữ liệu (Data Ingestion) và luồng truy vấn (Query Pipeline), được minh họa trong Hình \ref{fig:data-flow}:

\begin{figure}[H]
\centering
\begin{tikzpicture}[
    node distance=0.6cm and 0.3cm,
    box/.style={rectangle, draw, rounded corners, minimum height=1cm, minimum width=2.8cm, text centered, font=\small},
    data/.style={box, fill=blue!15, text width=2.6cm},
    process/.style={box, fill=orange!20, text width=2.6cm},
    storage/.style={box, fill=green!15, text width=2.6cm},
    user/.style={box, fill=red!12, text width=2.6cm},
    arrow/.style={->, thick, >=stealth},
    label/.style={font=\scriptsize\itshape, text=gray}
]

% Row 1: Data pipeline
\node[data] (raw) {Raw Data\\{\scriptsize arXiv JSON 4.8GB}};
\node[process, right=of raw] (clean) {Data Cleaning\\{\scriptsize Streaming + Filter}};
\node[data, right=of clean] (json) {Clean JSONL\\{\scriptsize 1.2M CS papers}};
\node[process, right=of json] (embed) {Embedding\\{\scriptsize GPU + Colab}};
\node[data, right=of embed] (vecjson) {JSONL + Vectors\\{\scriptsize 5.4GB, 384-dim}};

% Row 2: Indexing
\node[process, below=1.0cm of json] (bulk) {Bulk Ingestion\\{\scriptsize batch 5,000}};
\node[storage, right=of bulk] (estext) {arxiv\_text\\{\scriptsize BM25 Index 1.5GB}};
\node[storage, right=of estext] (esvec) {arxiv\_vectors\\{\scriptsize Vector Index 10.4GB}};

% Arrows - data pipeline
\draw[arrow] (raw) -- (clean);
\draw[arrow] (clean) -- (json);
\draw[arrow] (json) -- (embed);
\draw[arrow] (embed) -- (vecjson);

% Arrows - indexing
\draw[arrow] (vecjson.south) -- ++(0,-0.35) -| (bulk.north);
\draw[arrow] (bulk) -- (estext);
\draw[arrow] (estext) -- (esvec);

% Labels
\node[label, above=0.15cm of clean] {data\_prep\_full.py};
\node[label, above=0.15cm of embed] {embed\_data.py};
\node[label, above=0.15cm of bulk] {index\_data\_v2.py};

\end{tikzpicture}
\caption{Sơ đồ luồng nạp dữ liệu (Data Ingestion Pipeline) từ dữ liệu thô đến Elasticsearch}
\label{fig:data-flow}
\end{figure}

\section{Luồng Tìm kiếm Hybrid Search (Query Pipeline)}
Khi người dùng nhập truy vấn, hệ thống thực hiện tìm kiếm kết hợp (Hybrid Search) theo luồng được minh họa trong Hình \ref{fig:search-flow}:

\begin{figure}[H]
\centering
\begin{tikzpicture}[
    node distance=0.7cm and 0.4cm,
    box/.style={rectangle, draw, rounded corners, minimum height=0.9cm, minimum width=2.5cm, text centered, font=\small},
    user/.style={box, fill=red!12, text width=2.3cm},
    process/.style={box, fill=orange!20, text width=2.3cm},
    storage/.style={box, fill=green!15, text width=2.3cm},
    merge/.style={box, fill=violet!15, text width=2.3cm},
    arrow/.style={->, thick, >=stealth},
    darrow/.style={->, thick, >=stealth, dashed, blue!60}
]

% User input
\node[user] (query) {User Query\\{\scriptsize ``deep learning''}};
\node[process, right=of query] (api) {FastAPI Backend\\{\scriptsize parse + filter}};

% Two parallel branches
\node[process, above right=0.4cm and 1.5cm of api] (bm25q) {BM25 Search\\{\scriptsize title\^{}2, abstract}};
\node[process, below right=0.4cm and 1.5cm of api] (knnq) {Vector Search\\{\scriptsize cosine, k=10}};

% Encode step for kNN
\node[process, below=0.3cm of api] (encode) {Encode Query\\{\scriptsize MiniLM → 384d}};

% ES indices
\node[storage, right=of bm25q] (estext2) {arxiv\_text\\{\scriptsize Inverted Index}};
\node[storage, right=of knnq] (esvec2) {arxiv\_vectors\\{\scriptsize HNSW Graph}};

% RRF merge
\node[merge, right=3.5cm of api] (rrf) {RRF Merge\\{\scriptsize k=60, rank fusion}};

% Results
\node[user, right=of rrf] (result) {Top-10 Results\\{\scriptsize ranked by RRF}};

% Arrows
\draw[arrow] (query) -- (api);
\draw[arrow] (api) -- (bm25q);
\draw[arrow] (api) -- (encode);
\draw[arrow] (encode) -- (knnq);
\draw[arrow] (bm25q) -- (estext2);
\draw[arrow] (knnq) -- (esvec2);
\draw[darrow] (estext2.east) -- ++(0.3,0) |- (rrf.north);
\draw[darrow] (esvec2.east) -- ++(0.3,0) |- (rrf.south);
\draw[arrow] (rrf) -- (result);

\end{tikzpicture}
\caption{Luồng tìm kiếm Hybrid Search: BM25 và Vector chạy song song, gộp kết quả bằng RRF}
\label{fig:search-flow}
\end{figure}

\noindent Luồng truy vấn gồm các bước:
\begin{enumerate}
    \item Người dùng nhập truy vấn qua giao diện web. FastAPI Backend tiếp nhận, phân tích bộ lọc (năm, danh mục).
    \item \textbf{Luồng BM25:} Gửi truy vấn đến chỉ mục \texttt{arxiv\_text}, Elasticsearch phân tích cú pháp (tokenize, lowercase) và tính điểm BM25 dựa trên Inverted Index.
    \item \textbf{Luồng Vector search:} Mô hình \texttt{all-MiniLM-L6-v2} chuyển câu truy vấn thành vector 384 chiều. Vector này được gửi đến chỉ mục \texttt{arxiv\_vectors} để tìm $k$ láng giềng gần nhất trên đồ thị HNSW.
    \item \textbf{RRF Merge:} Hai danh sách kết quả được gộp bằng công thức RRF (xem Chương \ref{ch:search-implementation}). Các tài liệu xuất hiện ở thứ hạng cao trong cả hai luồng được đẩy lên đầu.
    \item Trả về top-10 kết quả cuối cùng cho giao diện người dùng kèm điểm RRF và metadata.
\end{enumerate}

\section{Kiến trúc hạ tầng (Infrastructure)}

\subsection{Công nghệ sử dụng}
\begin{itemize}
    \item \textbf{Python 3.x:} Xử lý dữ liệu (Data Processing) và tương tác với Elasticsearch thông qua thư viện \texttt{elasticsearch-py}.
    \item \textbf{Elasticsearch 8.12:} Search Engine cốt lõi, hỗ trợ Inverted Index, BM25 scoring, và filtering \cite{elasticsearch_docs}.
    \item \textbf{Docker \& Docker Compose:} Đóng gói và triển khai toàn bộ hạ tầng (Elasticsearch + Kibana) trong các container độc lập \cite{docker_docs}.
    \item \textbf{FastAPI:} Backend API cung cấp các endpoint RESTful cho giao diện tìm kiếm (xem Chương \ref{ch:ui}).
    \item \textbf{React + Vite:} Giao diện người dùng tương tác, kết nối FastAPI qua HTTP.
    \item \textbf{Kibana 8.12:} Giao diện quản trị và trực quan hóa dữ liệu, hỗ trợ Dev Tools để thử nghiệm truy vấn.
\end{itemize}

\subsection{Mô hình triển khai}
Hệ thống được triển khai dưới dạng Single-node cluster trên môi trường Docker cục bộ, sử dụng cấu hình Docker Compose như sau:

\begin{lstlisting}[caption={Docker Compose -- Elasticsearch \& Kibana},label={lst:docker-compose}]
services:
  elasticsearch:
    image: elasticsearch:8.12.2
    environment:
      - discovery.type=single-node
      - xpack.security.enabled=false
      - "ES_JAVA_OPTS=-Xms2g -Xmx2g"
    ports:
      - "9200:9200"
    volumes:
      - es_data:/usr/share/elasticsearch/data

  kibana:
    image: kibana:8.12.2
    environment:
      - ELASTICSEARCH_HOSTS=http://elasticsearch:9200
    ports:
      - "5601:5601"
\end{lstlisting}

\subsection{Đặc tính Big Data}
Mặc dù chạy trên single-node, kiến trúc bên trong của Elasticsearch vẫn duy trì các khái niệm phân tán như \textbf{Index}, \textbf{Shard} (phân mảnh dữ liệu) và \textbf{Replica} (bản sao). Điều này cho phép hệ thống dễ dàng mở rộng (scale-out) bằng cách thêm các node mới vào cluster khi khối lượng dữ liệu tăng lên mức hàng triệu bài báo mà không cần thay đổi kiến trúc ứng dụng.

\section{Index Mapping \& Bulk Strategy}

\subsection{Mapping Definition}
Lược đồ dữ liệu (Mapping) được định nghĩa tường minh để tối ưu hóa hiệu năng cho từng loại truy vấn:

\begin{lstlisting}[caption={Elasticsearch Index Mapping},label={lst:mapping}]
{
  "mappings": {
    "properties": {
      "id":         { "type": "keyword" },
      "title":      { "type": "text" },
      "abstract":   { "type": "text" },
      "categories": { "type": "keyword" },
      "authors":    { "type": "text" },
      "year":       { "type": "keyword" }
    }
  },
  "settings": {
    "number_of_shards": 1,
    "number_of_replicas": 0
  }
}
\end{lstlisting}

Trong đó:
\begin{itemize}
    \item Các trường \texttt{title}, \texttt{abstract}, \texttt{authors} được map dưới dạng \texttt{text} -- Elasticsearch tự động thực hiện phân tích cú pháp (tokenization, lowercasing, stop-word removal) và xây dựng Inverted Index để phục vụ tìm kiếm BM25.
    \item Các trường \texttt{id}, \texttt{year}, \texttt{categories} được map dưới dạng \texttt{keyword} -- lưu trữ nguyên bản, không qua phân tích, phục vụ exact-match filtering với tốc độ O(1).
\end{itemize}

\subsection{Bulk Indexing Strategy}
Thay vì gửi từng tài liệu qua HTTP POST riêng lẻ (rất chậm và tốn tài nguyên mạng), hệ thống sử dụng hàm \texttt{helpers.bulk()} của thư viện \texttt{elasticsearch-py}. Dữ liệu được gom thành các lô (batch) với kích thước \textbf{5,000 documents} mỗi lần gửi, giúp:
\begin{itemize}
    \item Giảm số lượng HTTP round-trip từ hàng triệu xuống còn vài trăm requests.
    \item Tối ưu hóa băng thông mạng và giảm tải cho bộ nhớ heap của Elasticsearch.
    \item Toàn bộ 1,203,108 documents được index trong thời gian hợp lý.
\end{itemize}

\section{Kiến trúc Dual-Index (Dual-Index Architecture)}
\label{sec:dual-index}

Để tối ưu hóa hiệu năng tìm kiếm và hỗ trợ nạp bài báo mới (incremental ingestion) mà không cần đánh chỉ mục lại toàn bộ, hệ thống được nâng cấp lên kiến trúc Dual-Index trên cùng một cụm Elasticsearch:

\begin{itemize}
    \item \textbf{Index \texttt{arxiv\_text}:} Chỉ chứa các trường văn bản (\texttt{title}, \texttt{abstract}, \texttt{authors}, \texttt{year}, \texttt{categories}), phục vụ tìm kiếm BM25. Kích thước chỉ 1.5GB cho 1.2 triệu bài báo do không lưu trữ vector nhúng.
    \item \textbf{Index \texttt{arxiv\_vectors}:} Chứa toàn bộ metadata và vector nhúng 384 chiều (\texttt{embedding}), phục vụ tìm kiếm vector. Kích thước 10.4GB.
\end{itemize}

\subsection{Cơ chế Nạp bài báo mới (Incremental Ingestion)}
Khi có bài báo mới, hệ thống tự động:
\begin{enumerate}
    \item Sinh vector nhúng 384 chiều bằng mô hình \texttt{all-MiniLM-L6-v2}.
    \item Đẩy metadata vào \texttt{arxiv\_text} qua Elasticsearch Index API.
    \item Đẩy ID và vector vào \texttt{arxiv\_vectors} qua Elasticsearch Index API.
    \item Refresh cả hai index để bài báo mới có thể tìm kiếm được ngay lập tức.
\end{enumerate}

Cơ chế này cho phép bổ sung bài báo mới trong thời gian thực mà không cần nạp lại toàn bộ 1.2 triệu bài báo từ đầu.